% !TEX program = xelatex
% ===========================================================
% Yanda Cheng -- Research CV (Medical Imaging / CV / ML / LLM) -- English-only Template
% Compile with: XeLaTeX (recommended)
% ===========================================================

\documentclass[letterpaper,10pt]{article}

% ---------- Packages ----------
\usepackage{iftex}
\usepackage[empty]{fullpage}
\usepackage{titlesec}
\usepackage[usenames,dvipsnames]{color}
\usepackage{enumitem}
\setlist[itemize]{noitemsep, topsep=3pt, left=1.5em}
\usepackage[hidelinks]{hyperref}
\usepackage{fancyhdr}
\usepackage{tabularx}

% ---------- Engine compatibility ----------
\ifPDFTeX
\fi
\ifXeTeX
  \usepackage{fontspec}
  \usepackage{xeCJK}
\fi
\ifLuaTeX
  \usepackage{fontspec}
  \usepackage{luatexja}
  \usepackage{luatexja-fontspec}
\fi

% ---------- Fonts ----------
\ifXeTeX
  \IfFontExistsTF{Microsoft YaHei}{
    \setCJKmainfont{Microsoft YaHei}
    \setCJKsansfont{Microsoft YaHei}
  }{
    \setCJKmainfont{Noto Sans CJK SC}
    \setCJKsansfont{Noto Sans CJK SC}
  }
\fi
\ifLuaTeX
  \IfFontExistsTF{Microsoft YaHei}{
    \setmainjfont{Microsoft YaHei}
  }{
    \setmainjfont{Noto Sans CJK SC}
  }
\fi
\linespread{1.03}\selectfont

% ---------- Page Style ----------
\pagestyle{fancy}
\fancyhf{}
\renewcommand{\headrulewidth}{0pt}
\renewcommand{\footrulewidth}{0pt}

% ---------- Margins ----------
\addtolength{\oddsidemargin}{-0.5in}
\addtolength{\textwidth}{1in}
\addtolength{\topmargin}{-0.9in}
\addtolength{\textheight}{1.15in}

% ---------- Section Formatting ----------
\titleformat{\section}{\large\bfseries}{}{0em}{}[\titlerule]
\titlespacing{\section}{0pt}{5pt plus 1pt minus 1pt}{3pt}

% ---------- Custom Commands ----------
\newcommand{\resumeItem}[1]{\item\small{#1}}

\newcommand{\resumeSubheading}[4]{
  \vspace{-1pt}\item
  \begin{tabular*}{0.97\textwidth}[t]{l@{\extracolsep{\fill}}r}
    \textbf{#1} & \small#2 \\
    \textit{\small#3} & \textit{\small#4} \\
  \end{tabular*}\vspace{-2pt}
}

\newcommand{\projheading}[2]{
  \vspace{1pt}\item
  \begin{tabular*}{0.97\textwidth}[t]{l@{\extracolsep{\fill}}r}
    \textbf{#1} & \textit{\small#2} \\
  \end{tabular*}\vspace{2pt}
}

\newcommand{\resumeSubHeadingList}{\begin{itemize}[leftmargin=0.0in, label={}]}
\newcommand{\resumeSubHeadingListEnd}{\end{itemize}}
\newcommand{\resumeItemListStart}{\begin{itemize}[leftmargin=0.2in]}
\newcommand{\resumeItemListEnd}{\end{itemize}}
\newcommand{\projItemListStart}{\begin{itemize}[leftmargin=0.2in, topsep=3pt]}
\newcommand{\projItemListEnd}{\end{itemize}\vspace{2pt}}

% =================================================================
\begin{document}
% =================================================================

% #################################################################
% ###################### ENGLISH VERSION ONLY #####################
% #################################################################

% ---------- Header ----------
\begin{center}
  {\huge \textbf{Yanda Cheng}}
\end{center}
\vspace{4pt}
\noindent\small
+1-859-338-3379 \quad$|$\quad
\href{mailto:nickcp39@gmail.com}{nickcp39@gmail.com} \quad$|$\quad
\href{https://yanda-cheng.com}{https://yanda-cheng.com} \quad$|$\quad
\href{https://www.linkedin.com/in/yanda-cheng-4888a216b}{LinkedIn} \quad$|$\quad
\href{https://scholar.google.com/citations?user=XkNfimMAAAAJ&hl=en}{Google Scholar}

\vspace{-4pt}

% ---------- Education ----------
\section{Education}

\begin{itemize}[leftmargin=0.0in, label={}, itemsep=2pt, topsep=2pt]
  \item \textbf{PhD} in Biomedical Engineering \hfill \textit{University at Buffalo (SUNY), Buffalo, NY} \hfill GPA: 4.0/4.0 \hfill 2021--2026 (Expected)
  \item \textbf{MEng} in Biomedical Engineering \hfill \textit{Cornell University (Ivy League), Ithaca, NY} \hfill 2020--2021
  \item \textbf{BS} in Electrical \& Computer Engineering \hfill \textit{University of Kentucky, Lexington, KY} \hfill 2015--2020
\end{itemize}

\vspace{-2pt}

% ---------- Work Experience ----------
\section{Work Experience}

\begin{itemize}[leftmargin=0in, label={}]

  % --- SGLang ---
  \vspace{-3pt}\item
  \begin{tabular*}{0.97\textwidth}[t]{l@{\extracolsep{\fill}}r}
    \textbf{SGLang (LLM/VLM Serving Framework)} & Mountain View, CA (remote) \\
    \textit{Open-Source Contributor} & \textit{Aug 2025 -- Feb 2026} \\
  \end{tabular*}\vspace{2pt}
  \begin{itemize}[leftmargin=0.2in, topsep=3pt]
    \item Added \textbf{Sequence Parallelism (SP)} support for \textbf{GLM-Image} in SGLang diffusion pipeline (\textbf{\#18077}): enabled multi-GPU sharded execution for DiT-style diffusion, established reproducible latency/VRAM baselines, and validated behavior across SP/TP settings to unblock high-resolution generation
    \item Fixed high-concurrency file-descriptor leaks in HTTP utils (\textbf{\#12047}) by ensuring urllib responses are always closed on success and error paths, preventing OSError (too many open files); stress-tested with 5,000 requests (100 threads) achieving 0 percent failures and p99 under 30 ms with stable FD count
    \item Unblocked out-of-the-box diffusion launches by fixing missing dependencies in \texttt{sglang[diffusion]} (\textbf{\#17900}; related \textbf{\#17671}): added \textbf{accelerate} and \textbf{ftfy} to extras, enabling \texttt{device\_map="cuda/auto"} model loading for Wan2.1 and FLUX-class models in clean-room installs
    \item Built and shared \textbf{Qwen} performance benchmarks focusing on \textbf{p99 latency} across multiple \textbf{SP/TP} configurations, providing actionable baselines for parallelism trade-offs and regression detection
  \end{itemize}

  \vspace{3pt}

  % --- Seno Medical ---
  \vspace{-3pt}\item
  \begin{tabular*}{0.97\textwidth}[t]{l@{\extracolsep{\fill}}r}
    \textbf{Seno Medical} & New York, NY \\
    \textit{Applied ML Engineer Intern} & \textit{May 2024 -- Aug 2024} \\
  \end{tabular*}\vspace{2pt}
  \begin{itemize}[leftmargin=0.2in, topsep=3pt]
    \item Collaborated with clinicians to curate 2,000+ medical imaging QA datasets with annotation guidelines; fine-tuned LLaMA (LoRA) and built a RAG system, cutting hallucination by 50 percent and raising Recall at 5 to 90 percent
    \item Designed an end-to-end evaluation harness with label schema (factuality, coverage, style), automatic metrics, and Python error-analysis scripts; reduced hallucination rate from 30 percent to 15 percent and increased Recall at 5 to 90 percent
    \item Implemented a modular RAG backend: chunking and ingestion jobs, FAISS-based vector store, retriever/re-ranker layer, and LLM orchestration, served via FastAPI and containerized with Docker for PHI-compliant on-prem deployment
    \item Set up Git-based workflows and observability (structured logs, latency/error-rate dashboards); shipped reproducible Docker images and deployment runbooks for downstream clinical teams
  \end{itemize}

  \vspace{3pt}

  % --- Roswell Park ---
  \vspace{-3pt}\item
  \begin{tabular*}{0.97\textwidth}[t]{l@{\extracolsep{\fill}}r}
    \textbf{Roswell Park Comprehensive Cancer Center} & Buffalo, NY \\
    \textit{Research Scientist} & \textit{May 2023 -- Aug 2023} \\
  \end{tabular*}\vspace{2pt}
  \begin{itemize}[leftmargin=0.2in, topsep=3pt]
    \item Designed a multi-modal data pipeline aligning imaging with structured text reports (de-identification, schema normalization, automated QC) to produce clean image--text pairs for generative modeling
    \item Built dataset generation jobs (filtering, sampling, augmentation, metadata tagging) to create balanced image--report datasets for conditional generation and downstream training
    \item Implemented a domain-specific latent diffusion model (Stable Diffusion--style, DDIM/DPMS sampling), trained on AWS, for conditional image-to-text report generation and synthetic data augmentation
    \item Automated GPU training and experiment tracking with standardized configs, logging, and comparison scripts, enabling fast iteration over model variants and dataset settings on the AWS stack
  \end{itemize}

  \vspace{3pt}

  % --- HydroSense ---
  \vspace{-3pt}\item
  \begin{tabular*}{0.97\textwidth}[t]{l@{\extracolsep{\fill}}r}
    \textbf{HydroSense Tech LLC} & Beijing / Singapore \\
    \textit{Co-founder \& Core Engineer} & \textit{2015 -- Present (part-time)} \\
  \end{tabular*}\vspace{2pt}
  \begin{itemize}[leftmargin=0.2in, topsep=3pt]
    \item Co-founded an IoT rainfall sensing startup and led ML for field-deployed devices, helping scale the product line from pilot deployments to \textbf{multi-million-dollar annual revenue} across municipal and industrial customers in 5+ countries
    \item Built an end-to-end deployment stack: C++ firmware on STM32-based controllers plus Python services on an edge gateway, with secure OTA updates over RS-485 (local bus) and Bluetooth/Wi-Fi (field apps), enabling \textbf{fleet-wide remote monitoring and configuration} for 5000+ devices
    \item Designed and shipped an \textbf{on-device 1D-CNN} for temperature-drift compensation using PyTorch training and INT8 quantization, reducing weight-estimation MAE by \textbf{15 percent in production} (20k+ installed sensors)
    \item Built an AI-driven calibration and prediction framework (Python/MLflow) to model nonlinearities and temperature dependence; reduced field drift by \textbf{30 percent} while tracking long-term device health across deployments
    \item \textbf{Patents:} Co-inventor on 3 patents, covering rainfall sensing hardware and ML-based drift compensation algorithms
  \end{itemize}

\end{itemize}

\vspace{-2pt}

% ---------- Research Projects ----------
\section{Selected Research Projects}

\begin{itemize}[leftmargin=0in, label={}]

  \projheading{Predictive modeling of chronic foot ulcer outcomes using longitudinal photoacoustic imaging}{npj Imaging, 2026}
  \projItemListStart
    \item Designed longitudinal follow-up protocol for over 400 diabetic foot ulcer patients; coordinated \textbf{IRB approval}, multi-center patient enrollment, and data collection with endocrinology, vascular surgery, and wound care teams
    \item Combined radiomics features with LASSO/Random Forest models achieving AUC over 0.85 for quantitative ulcer healing prediction, directly supporting clinical decision-making
    \item \textit{Applicable to:} longitudinal outcome prediction across medical imaging modalities, risk stratification in clinical cohorts, and time-series modeling in CV/ML pipelines
  \projItemListEnd

  \projheading{OneTouch automated photoacoustic and ultrasound imaging of breast in standing pose}{IEEE Trans.\ Medical Imaging, 2025, cited 8 times}
  \projItemListStart
    \item Co-developed a multi-modal clinical imaging platform integrating AI-assisted diagnosis for early breast cancer screening; coordinated with radiologists and surgeons on clinical protocol design and image interpretation validation
    \item \textit{Applicable to:} multi-modal medical imaging, computer vision for cancer screening, and deployment of AI tools in radiology workflows
  \projItemListEnd

  \projheading{Dual-scan photoacoustic tomography for the imaging of vascular structure on foot}{IEEE Trans.\ UFFC, 2023, cited 21 times}
  \projItemListStart
    \item Developed quantitative 3D vascular imaging protocol for peripheral arterial disease (PAD) assessment; established clinical imaging standards in collaboration with a vascular surgery team
    \item \textit{Applicable to:} quantitative vascular imaging, angiography analysis, and 3D vessel segmentation in medical imaging and CV applications
  \projItemListEnd

\end{itemize}

\vspace{-2pt}

% ---------- Publications ----------
\section{Selected Publications (First / Co-Author, Peer-Reviewed)}

\begin{itemize}[leftmargin=0.2in, itemsep=1pt]
  \item \small \textbf{Cheng Y.}, Huang C., et al. Predictive modeling of chronic foot ulcer outcomes using longitudinal photoacoustic imaging. \textit{npj Imaging} 4(1), 12, \textbf{2026}.
  \item \small \textbf{Cheng Y.}, Huang C., Bing R.W., et al. Dysphagia assessment based on photoacoustic imaging: a pilot ex vivo and in vivo study in infant swine. \textit{Med-X} 3(1), 1--9, \textbf{2025}.
  \item \small \textbf{Cheng Y.}, Zheng W., et al. Unsupervised denoising of photoacoustic images based on the Noise2Noise network. \textit{Biomed.\ Optics Express} 15(8), 4390--4405, \textbf{2024}.\quad(14 citations)
  \item \small \textbf{Cheng Y.}, Huang C., et al. Dual-scan photoacoustic tomography for the imaging of vascular structure on foot. \textit{IEEE Trans.\ UFFC} 70, \textbf{2023}.\quad(21 citations)
  \item \small Huang C., \textbf{Cheng Y.}, et al. Enhanced clinical photoacoustic vascular imaging through skin localization and adaptive weighting. \textit{Photoacoustics} 42, 100690, \textbf{2025}. [4 citations]
  \item \small Zhang H., ..., \textbf{Cheng Y.}, et al. OneTouch automated PA/US breast imaging in standing pose. \textit{IEEE Trans.\ Medical Imaging}, \textbf{2025}. (8 citations)
  \item \small Liu X., \textbf{Cheng Y.}, et al. Simultaneous tissue blood flow and oxygenation with a wearable fiber-free optical sensor. \textit{J.\ Biomed.\ Optics} 26(1), 012705, \textbf{2021}.\quad\textbf{[38 citations]}
\end{itemize}
\vspace{2pt}
\noindent\small\textbf{Total: about 100 citations (Google Scholar); 11 SCI papers; Venues: npj Imaging, IEEE TMI, IEEE TUFFC, Photoacoustics, BOE, Med-X, JBO}

\vspace{-2pt}

% ---------- Technical Skills ----------
\section{Technical Skills}
\noindent
\small
\textbf{Medical Imaging:} Image reconstruction, denoising (DL/classical), registration, segmentation, 3D visualization; CT, PA, US, multiphoton microscopy \\
\textbf{AI / Machine Learning / LLMs:} CNN, UNet, Transformer, LSTM, Radiomics, LASSO, SVM, Random Forest, deep learning denoising; large language models, RAG systems, prompt engineering, multimodal models \\
\textbf{Clinical Data Analysis:} Longitudinal analysis, Kaplan-Meier survival analysis, ROC/AUC, regression, statistical testing (Python\,/\,R\,/\,SPSS) \\
\textbf{Research Workflow:} IRB/ethics protocol submission, multi-site data management, SCI manuscript preparation \& submission, academic presentation \\
\textbf{Programming:} Python, C/C++, MATLAB, R, SQL; PyTorch, scikit-learn, OpenCV, Pandas, Docker \\
\textbf{Languages:} Mandarin Chinese (native), English (fluent -- academic writing, cross-national team collaboration)

\vspace{-2pt}

% ---------- Awards ----------
\section{Awards \& Honors}
\begin{itemize}[leftmargin=0.2in, itemsep=1pt]
  \item \small \textbf{\$3,000 Startup Seed Grant}, University at Buffalo Entrepreneurship Award -- PA-based Breast Cancer Detection (2023)
  \item \small Key Contributor, NIH-Funded Research Project (PI: Prof.\ Jun Xia, University at Buffalo; ongoing)
  \item \small Outstanding Teaching Assistant -- BME 503 Image Processing \& BME 302 Medical Devices, University at Buffalo
\end{itemize}

\end{document}

